\documentclass{beamer}
\usepackage[english,greek]{babel}
\usecolortheme{default}
\usetheme{Madrid}
\usepackage{tikz}
\usetikzlibrary{positioning} % Βιβλιοθήκη εντολών και δομών θέσης. 
\usetikzlibrary{automata} % Βιβλιοθήκη αυτομάτων 
\usetikzlibrary{shapes} 
\usetikzlibrary {arrows.meta}
\usepackage{amsmath}
\usepackage{colortbl}
\usepackage{cancel}

\AtBeginSection[]
{
	\begin{frame}
		\frametitle{Πίνακας Περιεχομένων}
		\tableofcontents[currentsection]
	\end{frame}
}
%Information to be included in the title page:
\title[ΨΕΣ - Έξτρα Υλικό] %optional
{Ψηφιακή Επεξεργασία Σήματος - Έξτρα Υλικό 01}

\subtitle{Τριγωνομετρία \& Άσκηση Δειγματοληψίας (και λοιπά κουραφέξαλα)}

\author[\textlatin{bruhvius23}] % (optional, for multiple authors)
{\textlatin{bruhvius23}\inst{1}}

\institute[ΠΑΔΑ] % (optional)
{
	\inst{1}%
	Τμήμα Μηχανικών Πληροφορικής και Υπολογιστών\\
	Πολυτεχνική Σχολή (χαχα καλό ε) \\
	Πανεπιστήμιο Δυτικής Αττικής
}

\date[Γ\textlatin{WU} 2026] % (optional)
{Γεωργουλάκη \textlatin{Weekly Update}, Οκτώβριος 2026}

\begin{document}

	\frame{\titlepage}
% -----------------
\section{Ασκήσεις Δειγματοληψίας + Τριγωνομετρία}	
% δειγματοληψία
% ---------------	
	\begin{frame}
		\frametitle{Ασκήσεις Δειγματοληψίας - 1η (Εκφώνηση)}
		\textbf{Παράδειγμα:} Δίνεται το αναλογικό σήμα $x(t) = 2cos(100\pi t)$. Θεωρώντας ότι λαμβάνονται δείγματα του σήματος με ρυθμό $F_s=75Hz$, ποιο το σήμα διακριτού χρόνου που προκύπτει μετά τη δειγματοληψία και ποιο το σήμα συνεχούς χρόνου που προκύπτει από την ανακατασκευή$;$
	\end{frame}
	
	\begin{frame}
	\frametitle{Ασκήσεις Δειγματοληψίας - 1η (Λύση-Εξήγηση)}
	Σημειώνουμε αρχικά το ότι έχουμε ένα ημιτονοειδές σήμα. Αυτά, μπορούν να εκφραστούν ή ως $Acos(\Omega t)$ (ή $Asin(\Omega t)$) , με A να ορίζεται το πλάτος του σήματος και ω η \textbf{κυκλική (ή γωνιακή) συχνότητα}, ή ως $Acos(2\pi F t)$ και αντίστοιχα για το ημίτονο $sin$, όπου $\Omega = 2 \pi F$. 
	\end{frame}
	
	\begin{frame}
			\frametitle{Παράδειγμα Αναλογικού Σήματος - Ημίτονο}
		Έστω η γραφική παράσταση του κλασικού ημιτόνου $sin(x)$. \pause
		\begin{center}
		\begin{tikzpicture}
			\draw[->] (-2.5,0) -- (7.5,0);
			\draw[->] (0,-2.5) -- (0,2.5);
			\draw (0,0) sin (1.57079632679,1) cos(3.1415926535,0) sin(4.71238898038,-1) cos(6.28318,0);
			% Σημείωση 0 -----
			\draw (0,0) node[below left,shape=circle] {0};
			% -----------
			
			%---------
			% Σημείωση π/2
			% π/2 στον x
			\draw[blue,thick,dotted] (1.57079632679,1) -- (1.57079632679,0); 
			\draw (1.57079632679,0) node[below] {$\frac{\pi}{2}$};
			% 1 στον y
			\draw[blue,thick,dotted] (1.57079632679,1) -- (0,1);
			\draw (0,1) node[left] {$1$};
			% circle, fill, inner sep=1pt
			\draw (1.57079632679,1) node[circle, fill, inner sep=1pt] {};
			% ---------
			
			% Σημείωση π
			% π στον x

			% 0 στον y

			% circle, fill, inner sep=1pt
			\draw (3.1415926535,0) node[circle, fill, inner sep=1pt]{};			
			\draw (3.1415926535,0) node[below left] {$\pi$};
			
		 	%---------
		 	% Σημείωση 3π/2
		 	% 3π/2 στον x
		 	\draw[blue,thick,dotted] (4.71238898038,-1) -- (4.71238898038,0); 
		 	\draw (4.71238898038,0) node[above] {$\frac{3\pi}{2}$};
		 	% 1 στον y
		 	\draw[blue,thick,dotted] (4.71238898038,-1) -- (0,-1);
		 	\draw (0,-1) node[left] {$-1$};
		 	% circle, fill, inner sep=1pt
		 	\draw (4.71238898038,-1) node[circle, fill, inner sep=1pt] {};
		 	% ---------
			
			
			% Σημείωση 2π
			% 2π στον x
			
			% 0 στον y
			
			% circle, fill, inner sep=1pt
			\draw (6.28318,0) node[circle, fill, inner sep=1pt]{};			
			\draw (6.28318,0) node[below right] {$2\pi$};
		\end{tikzpicture}
		\end{center} \pause
		Παρατηρήστε ότι το σήμα φτάνει μέχρι το +1 και το -1 στον κατακόρυφο άξονα. Αν λάβουμε υπόψη την αρχή των αξόνων $O$, απέχει 1 μονάδα από την μέγιστη και ελάχιστη τιμή του σήματος αντίστοιχα. Αυτό καθιστά το πλάτος του $A$, όπου εδώ ισχύει $A=1$.
	\end{frame}

	\begin{frame}
	\frametitle{Παράδειγμα Αναλογικού Σήματος - Ημίτονο}
	Έστω η γραφική παράσταση του κλασικού ημιτόνου $sin(x)$. 
	\begin{center}
		\begin{tikzpicture}
	\draw[->] (-2.5,0) -- (7.5,0);
	\draw[->] (0,-2.5) -- (0,2.5);
	\draw (0,0) sin (1.57079632679,1) cos(3.1415926535,0) sin(4.71238898038,-1) cos(6.28318,0);
	% Σημείωση 0 -----
	\draw (0,0) node[below left,shape=circle] {0};
	% -----------
	
	%---------
	% Σημείωση π/2
	% π/2 στον x
	\draw[blue,thick,dotted] (1.57079632679,1) -- (1.57079632679,0); 
	\draw (1.57079632679,0) node[below] {$\frac{\pi}{2}$};
	% 1 στον y
	\draw[blue,thick,dotted] (1.57079632679,1) -- (0,1);
	\draw (0,1) node[left] {$1$};
	% circle, fill, inner sep=1pt
	\draw (1.57079632679,1) node[circle, fill, inner sep=1pt] {};
	% ---------
	
	% Σημείωση π
	% π στον x
	
	% 0 στον y
	
	% circle, fill, inner sep=1pt
	\draw (3.1415926535,0) node[circle, fill, inner sep=1pt]{};			
	\draw (3.1415926535,0) node[below left] {$\pi$};
	
	%---------
	% Σημείωση 3π/2
	% 3π/2 στον x
	\draw[blue,thick,dotted] (4.71238898038,-1) -- (4.71238898038,0); 
	\draw (4.71238898038,0) node[above] {$\frac{3\pi}{2}$};
	% 1 στον y
	\draw[blue,thick,dotted] (4.71238898038,-1) -- (0,-1);
	\draw (0,-1) node[left] {$-1$};
	% circle, fill, inner sep=1pt
	\draw (4.71238898038,-1) node[circle, fill, inner sep=1pt] {};
	% ---------
	
	
	% Σημείωση 2π
	% 2π στον x
	
	% 0 στον y
	
	% circle, fill, inner sep=1pt
	\draw (6.28318,0) node[circle, fill, inner sep=1pt]{};			
	\draw (6.28318,0) node[below right] {$2\pi$};
	\end{tikzpicture}
	\end{center}
	Υποθέστε επίσης πως το σύστημα των συντεταγμένων εννοεί τα $y$ και $x$ στον κατακόρυφο και οριζόντιο άξονα αντίστοιχα αλλά δεν αναγράφονται ρητά για εξοικονόμηση χώρου.
\end{frame}	

	\begin{frame}
	\frametitle{Παράδειγμα Αναλογικού Σήματος - Ημίτονο}
	\begin{center}
		\begin{tikzpicture}
			\draw[->] (-2.5,0) -- (7.5,0);
			\draw[->] (0,-2.5) -- (0,2.5);
			\draw (0,0) sin (1.57079632679,1) cos(3.1415926535,0) sin(4.71238898038,-1) cos(6.28318,0);
			% Σημείωση 0 -----
			\draw (0,0) node[below left,shape=circle] {0};
			% -----------
			
			%---------
			% Σημείωση π/2
			% π/2 στον x
			\draw[blue,thick,dotted] (1.57079632679,1) -- (1.57079632679,0); 
			\draw (1.57079632679,0) node[below] {$\frac{\pi}{2}$};
			% 1 στον y
			\draw[blue,thick,dotted] (1.57079632679,1) -- (0,1);
			\draw (0,1) node[left] {$1$};
			% circle, fill, inner sep=1pt
			\draw (1.57079632679,1) node[circle, fill, inner sep=1pt] {};
			% ---------
			
			% Σημείωση π
			% π στον x
			
			% 0 στον y
			
			% circle, fill, inner sep=1pt
			\draw (3.1415926535,0) node[circle, fill, inner sep=1pt]{};			
			\draw (3.1415926535,0) node[below left] {$\pi$};
			
			%---------
			% Σημείωση 3π/2
			% 3π/2 στον x
			\draw[blue,thick,dotted] (4.71238898038,-1) -- (4.71238898038,0); 
			\draw (4.71238898038,0) node[above] {$\frac{3\pi}{2}$};
			% 1 στον y
			\draw[blue,thick,dotted] (4.71238898038,-1) -- (0,-1);
			\draw (0,-1) node[left] {$-1$};
			% circle, fill, inner sep=1pt
			\draw (4.71238898038,-1) node[circle, fill, inner sep=1pt] {};
			% ---------
			
			
			% Σημείωση 2π
			% 2π στον x
			
			% 0 στον y
			
			% circle, fill, inner sep=1pt
			\draw (6.28318,0) node[circle, fill, inner sep=1pt]{};			
			\draw (6.28318,0) node[below right] {$2\pi$};
		\end{tikzpicture}
	\end{center}
	Σημειώστε τώρα πως έχουμε τη μορφή $Asin(\Omega x)$ , δεν φαίνεται έστω κάποιο 2π. Εδώ, θα πρέπει για να βρούμε τα μεγέθη που θέλουμε, να αναφερθούμε σε κάποιους παλιούς ορισμούς.
\end{frame}
\begin{frame}
	\frametitle{Παράδειγμα Αναλογικού Σήματος - Ημίτονο (Ορισμοί)}
	Αν βρεθούμε στο σημείο $(2\pi,0)$ και προβούμε στο σημείο όπου $x=0$, τότε έχουμε $sin(x)=sin(x-2\pi)=0$. Ομοίως και αν προσθέσουμε $2\pi$ όταν $x=0$ θα βρεθούμε στο ίδιο σημείο. Βάσει τριγωνομετρίας κιόλας όπου θα δούμε πρόχειρα πως εξηγείται αυτό, αν συνεχίζαμε κατά 2π απ'όπου $x=2\pi$, τότε θα καταληγάμε με $sin(x)=sin(x-2\pi)=sin(x+2\pi)=0$. \pause
	\\~\\
	Προσπαθήστε να κάνετε το ίδιο για τα σημεία όπου ξέρετε ότι $sin(x) = 1$ ή που ξέρετε την τιμή της συνάρτησης, να δείτε αν $sin(x-2\pi) = sin(x+2\pi) = sin(x)$. Θα προσέξετε το εξής: 	Θα προσέξετε ότι το σήμα \textbf{επαναλαμβάνεται} άνα $2\pi$ χρονικές μονάδες. Αυτό ονομάζεται \textbf{περιόδος} ή \textbf{$T$}. Εδώ, αφού επαναλαμβάνεται το σήμα ανα $2\pi$ χρονικές μονάδες έχουμε $T=2\pi$ (έστω $T=2\pi ~sec$).
\end{frame}
\begin{frame}
	\frametitle{Παράδειγμα Αναλογικού Σήματος - Ημίτονο (Ορισμοί)}

	
	\begin{block}{Περίοδος - Ορισμός}
	Η \textbf{περίοδος}, συμβολιζόμενη ως $T$, αποτελεί τον ελάχιστο χρόνο μεταξύ δύο επαναλήψεων ενός περιοδικού σήματος. Μετριέται συνήθως σε $sec$.
	\end{block}
	Η περίοδος και το αν ένα σήμα είναι περιοδικό επικυρώνονται με δύο συνθήκες: \pause 
	\begin{itemize} 
		\item $\exists T>0,\forall x \in A: x-T \in A \wedge x+T \in A$, όπου $A$ το πεδίο ορισμού της συνάρτησης (ή γενικά με όσα $x$ μπορούμε να ``παίξουμε μπάλα``). \pause
		\item Μετά ή πρίν από Τ χρονικές μονάδες από ένα δοθέν σημείο, η συνάρτηση να επιστρέφει την ίδια τιμή δηλαδή $f(x-T) = f(x+T) = f(x)$.
	\end{itemize}
\end{frame}

\begin{frame}
	\frametitle{Παράδειγμα Αναλογικού Σήματος - Ημίτονο (Ορισμοί)}
	Ωστόσο, βλέπουμε πως η μορφή του ημιτονοειδούς σήματος που έχουμε ($sin(x)$), δεν αναγράφει κάπου την περιόδο του ή κατι κοντινό σε αυτή. Εδώ, δεδομένου ότι δεν μιλάμε \textbf{απλώς} για την συνάρτηση ημιτόνου, μπορούμε να χρησιμοποιήσουμε έναν απλό (αλλά προσωρινά άχρηστο) ορισμό, προιδεασμένοι από την σχέση $\Omega=2\pi F$:
	\begin{center}
		Η κυκλική συχνότητα $\Omega$ αφορά στο ρυθμό μεταβολής ενός γωνιακού μεγέθους και \alert{συμπεριλαμβάνει} την συχνότητα που θέλουμε να βρούμε.
	\end{center}
	Η συχνότητα, με τη σειρά της, ορίζεται ως εξής:
		\begin{block}{Συχνότητα - Ορισμός}
		Η \textbf{συχνότητα}, συμβολιζόμενη ως $F$, αποτελεί τον αριθμό των επαναλήψεων ενός περιοδικού σήματος εντός μίας χρονικής μονάδας (π.χ. $sec$). Μετριέται συνήθως σε \textlatin{$Hz$ (Hertz)} ή \textlatin{sec$^{-1}$}.
	\end{block}
\end{frame}

\begin{frame}
	\frametitle{Παράδειγμα Αναλογικού Σήματος - Ημίτονο (Ορισμοί)}
	Η μονάδα μέτρησης της συχνότητας προκύπτει από το γεγονός ότι η περίοδος και η συχνότητα είναι \textbf{αντίστροφα} μεγέθη (δηλ. $F=\frac{1}{T}$).\\
	Εδώ όμως, βρίσκουμε μόνο την γωνιακή συχνότητα $\Omega$, όπου αν βασιστούμε στη συνάρτηση $sin(x)$ , έχουμε ότι $\Omega = 1$ αν υποθέσουμε τη μορφή $Asin(\Omega t)$.
	\\ 
	Για να βρούμε τη συχνότητα, απλώς κάνουμε αντικαταστάσεις.
\end{frame}

\begin{frame}
	\frametitle{Παράδειγμα Αναλογικού Σήματος - Ημίτονο (Ορισμοί)}
	Έστω $\Omega = 1$ και $\Omega = 2 \pi F$. Εκ των δύο σχέσων απορέει: \pause
	\begin{center}
		\begin{align*}
			\Omega = 1 \rightarrow 2\pi F = 1 \rightarrow F = \frac{1}{2\pi}  
		\end{align*}
	\end{center}
	Άρα η συχνότητα του $sin(x)$ είναι $\frac{1}{2\pi}$ (έστω $Hz$), πράγμα το οποίο μπορούμε να επαληθεύσουμε και με τον τύπο $F=T^{-1}$. Ωστόσο, θα χρησιμοποιήσουμε αυτόν τον τύπο ώστε να αντλήσουμε την συχνότητα που θέλουμε να εξετάσουμε.
\end{frame}

\begin{frame}
	\frametitle{Ασκήσεις Δειγματοληψίας - 1η (Λύση-Εξήγηση)}
	Γυρνώντας στο σήμα της άσκησης, έχουμε $x(t) = 2\cos(100\pi t)$.  \\~\\ Εδώ, μπορούμε με ευκολία να υποθέσουμε πως έχουμε ένα σήμα της μορφής $A\cos(\Omega t)$, όπου $A = 2$ και $\Omega = 100\pi$. \\~\\
	\pause
	
	Γνωρίζοντας ότι $\Omega = 2\pi F$, αντικαθιστούμε τις μεταβλητές με γνωστές τιμές στην προηγούμενη σχέση, καταλήγοντας στην $100\pi = 2\pi F \rightarrow \frac{100\pi}{2\pi} = F \rightarrow F = 50Hz$.
\end{frame}

\begin{frame}
		\frametitle{Ασκήσεις Δειγματοληψίας - Δειγματοληψία}
		Πριν συνεχίσουμε, ας σημειώσουμε στην... κοινή Ελληνική το τι πρόκειται να δούμε. \\~\\
		Όταν λαμβάνουμε ένα σήμα αναλογικό (συνεχούς χρόνου και συνεχούς πλάτους), αυτό σημαίνει ότι έχουμε να κάνουμε με ένα σήμα \alert{\textbf{άπειρων}} υποδιαιρέσεων στον \textbf{χρόνο} (ακόμη και σε ένα διάστημα 2```` μπορεί να έχουμε 0.1 δεύτερα , 0.01, 0.001 κ.ο.κ.) , με το ίδιο πράγμα να ισχύει και για το \textbf{πλάτος} (παρόμοια λογική με το ότι μεταξύ του 0 και του 1 υπάρχουν άπειροι αριθμοί).
\end{frame}

\begin{frame}
	\frametitle{Ασκήσεις Δειγματοληψίας - Δειγματοληψία}
	Αυτό, όπως αντιλαμβάνεστε, θα ζορίσει πάρα πολύ τον υπολογιστή η οποιοδήποτε σύστημα κληθεί να αποθηκεύσει αυτό το σήμα γιατί θα έχει υπέρογκα δεδομένα να διαχειριστεί! \\~\\
	Έτσι, κληθήκαμε τώρα να αντιμετωπίσουμε δύο θέματα:
	\begin{itemize}
		\item 	\textbf{Πως} θα διαχειριστούμε τον συνεχή χρόνο \pause
		\item \textbf{Πως} θα διαχειριστούμε το συνεχές πλάτος
	\end{itemize}

\end{frame}

\begin{frame}
	\frametitle{Ασκήσεις Δειγματοληψίας - Δειγματοληψία}
	Για αυτό το κομμάτι ενισχυτικού υλικού, θα ασχοληθούμε με το πρώτο θέμα, για το οποίο έχει υπάρξει μέριμνα να αντιμετωπιστεί από τα ψηφιακά συστήματα. Πως όμως? \pause 
	\\~\\ 
	Μια ιδέα είναι να σκεφτούμε το εξής: τι και αν αντί για κάθε χρονική στιγμή του σήματος... αποθηκεύαμε ανά συγκεκριμένα τακτά διαστήματα τις τιμές του, σαν να ``τραβάμε ένα \textlatin{screenshot}'' του σήματος?
	\pause \\~\\
	Αυτή η ιδέα, όσο απίστευτο και να φαντάζει, είναι αυτή που εφαρμόζεται γενικά και αποτελεί το πρώτο εκ των κύριων τριών σταδίων της μετατροπής αναλογικού σήματος σε ψηφιακό. Αυτή είναι, λοιπόν, η λεγόμενη \textbf{δειγματοληψία} (\textlatin{sampling}) η οποία ονομάζεται έτσι καθώς λαμβάνουμε δείγματα του σήματος ανά τα προαναφερθέντα διαστήματα.
	
\end{frame}
\begin{frame}
	\frametitle{Ασκήσεις Δειγματοληψίας - Δειγματοληψία}
	Πριν συνεχίσουμε, οφείλουμε να σημειώσουμε δύο ορισμούς που θα μας φανούν εξαιρετικά χρήσιμοι στη συνέχεια.
	\begin{block}{Περίοδος Δειγματοληψίας - Ορισμός}
		Η \textbf{περίοδος δειγματοληψίας} (\textlatin{sampling period}) αποτελεί το χρονικό διάστημα μεταξύ της λήψης 2 διαδοχικών δειγμάτων. Συμβολίζεται με $T_s$.
	\end{block}
		\begin{block}{Συχνότητα Δειγματοληψίας - Ορισμός}
		Η \textbf{συχνότητα (ή ρυθμός) δειγματοληψίας} (\textlatin{sampling frequency / sampling rate}) αποτελεί τον αριθμό των δειγμάτων που λαμβάνονται σε μια χρονική μονάδα (π.χ. δευτερόλεπτα). Συμβολίζεται με $F_s$.
	\end{block}
	Για παράδειγμα, αν παίρνουμε ένα δείγμα ανά $0.5\sec$ τοτε $T_s=0.5\sec$ και $F_s = 1/0.5 Hz \rightarrow F = 2 Hz$.
\end{frame}


\begin{frame}
		\frametitle{Ασκήσεις Δειγματοληψίας - Δειγματοληψία (συν.)}
	Γνωρίζουμε λοιπόν ότι η συχνότητα του αναλογικού μας σήματος είναι $50Hz$. Η δειγματοληψία σαν τεχνική διέπεται από ένα συγκεκριμένο θεώρημα το οποίο άπαξ και αληθεύει, καθιστά δυνατή την ορθή ανακατασκευή ενός δειγματοληπτημένου αναλογικού σήματος (διακριτού χρόνου) στο αντίστοιχο αναλογικό σήμα. 
	\begin{block}{Θεώρημα \textlatin{Nyquist-Shannon}}
		Για να μπορεί να ανασυσταθεί ένα αναλογικό σήμα από ένα ψηφιακό δίχως απώλεια πληροφορίας, θα πρέπει $F_s \geq 2F_{max}$.
	\end{block}
	\pause
	Το $F_{max}$ συνιστά τη μέγιστη συχνότητα ενός σήματος, που συνήθως θα είναι αυτή που θα εκλαμβάνεται από το εκάστοτε σήμα (π.χ. για το παρόν σήμα έχουμε $F_{max} = 50Hz$ καθώς $F=50Hz$).
\end{frame}

\begin{frame}
		\frametitle{Ασκήσεις Δειγματοληψίας - 1η (Λύση-Εξήγηση)}
		Εδώ, βλέπουμε πως $F_s < 2*F_{max} \rightarrow 75 < 2*50 \rightarrow 75 < 100$. Επομένως, \textbf{ΔΕΝ} ικανοποιείται η συνθήκη του θεωρήματος \textlatin{Nyquist} και θα αντιμετωπίσουμε θέμα. \pause \\~\\
		Συγκεκριμένα, μπορεί να καταλήξουμε να χειριζόμαστε και κάποιες συχνότητες που κανονικά δεν υπάρχουν (ψευδείς συχνότητες). Εδώ δεν πρόκειται να σημειωθεί κάτι για την απόδειξη-εξήγηση, καθώς αυτή αναγράφεται στις σημειώσεις της 1ης εβδομάδας. \\~\\
		Για τώρα, ας επικεντρωθούμε στην κατασκευή του σήματος διακριτού χρόνου που αποτελεί επακόλουθο της (κακής) δειγματοληψίας μας.
\end{frame}

\begin{frame}
	\frametitle{Ασκήσεις Δειγματοληψίας - 1η (Λύση-Εξήγηση)}
	Αρχικά, θα πάρουμε το αναλογικό μας σήμα και την μορφή που γνωρίζουμε ($A\cos(\Omega t)$) και θα αντικαταστήσουμε το $\Omega t$ με το $\omega n$. 
	\begin{itemize}
		\item Το $\omega$ αποτελεί την ψηφιακή κυκλική συχνότητα, η οποία ορίζεται ως $\omega = \Omega T_s = \frac{\Omega}{F_s}$.
		\item Το $n$ αποτελεί έναν απαριθμητή των δειγμάτων δηλ. μια μεταβλητή που σημείωνε ότι βρισκόμαστε στο 1ο,2ο, ... $n$-οστό δείγμα.
	\end{itemize} 
	Στην παρούσα περίπτωση έχουμε $x(t)=2\cos(100\pi t) \rightarrow x[n] = 2\cos(\frac{100\pi}{75}n) \rightarrow x[n] = 2\cos(\frac{4\pi n}{3})$ (διαιρώντας τον αριθμητή και τον παρονομαστή με 25).
\end{frame}


\begin{frame}
	\frametitle{Ασκήσεις Δειγματοληψίας - Τριγωνομετρία}
	Τώρα, το θέμα είναι να καταλήξουμε με μια πιο απλή μορφή του συνημιτόνου (προτιμότερα, κάποια με μικρότερη συχνότητα). Πώς θα το καταφέρουμε αυτό$\colon$
	\pause \\~\\
	Για να δούμε πως θα το επιτύχουμε αυτό, ας αναφερθούμε στον τριγωνομετρικό κύκλο!
\end{frame}
\begin{frame}
		\frametitle{Ασκήσεις Δειγματοληψίας -  Τριγωνομετρία}
		\begin{center}
			\begin{tikzpicture}
		\draw (-3,0) -- (3,0);
		\draw (0,3) -- (0,-3);
		\draw (0,0) circle (2);
		\draw (3,0) node[right] {$x$};
		\draw (-3,0) node[left] {$-x$};
		\draw (0,3) node[right] {$y$};
		\draw (0,-3) node[right] {$-y$};
	\end{tikzpicture}		
		\end{center}

\end{frame}


\begin{frame}
	\frametitle{Ασκήσεις Δειγματοληψίας -  Τριγωνομετρία}
	\begin{columns}
			\column{0.4\textwidth}

			\begin{tikzpicture}
				\draw (-2,0) -- (2,0);
				\draw (0,2) -- (0,-2);
				\draw (0,0) circle (1);
				\draw (2,0) node[right] {$x$};
				\draw (-2,0) node[left] {$-x$};
				\draw (0,2) node[right] {$y$};
				\draw (0,-2) node[right] {$-y$};
			\end{tikzpicture}		

		\column{0.6\textwidth}
		\centering
		Εδώ, σε μικρότερη κλίμακα, είναι ο τριγωνομετρικός κύκλος. \\~\\
		Αυτός αποτελεί μια ειδική περίπτωση κύκλου, όπου η ακτίνα του δηλ. η απόσταση του κέντρου του από την περιφέρεια του έχει μήκος \textbf{μονάδα} ($r=1$).
	\end{columns}
	
\end{frame}

\begin{frame}
	\frametitle{Ασκήσεις Δειγματοληψίας -  Τριγωνομετρία}
	\begin{columns} 
		\column{0.6\textwidth}
		\begin{tikzpicture}
			\draw (-3,0) -- (3,0);
			\draw (0,3) -- (0,-3);
			\draw (0,0) circle (2);
			\draw (3,0) node[right] {$x$};
			\draw (-3,0) node[left] {$-x$};
			\draw (0,3) node[right] {$y$};
			\draw (0,-3) node[right] {$-y$};
			% radius and triangle
			\draw[red] (0,0) -- (1.2,1.6);
			\draw[thick, blue] (1.2,1.6) -- (1.2,0);
			\draw[thick,orange] (0,0) -- (1.2,0);
			\draw (1.0,1.5) node[below left] {{\textcolor{red}{$r$}}};
			\draw (0.6,-0.15) node {{\textcolor{orange}{$x'$}}};
			\draw (1.20,1.20) node[below right] {{\textcolor{blue}{$y'$}}};
			%
			 \draw (0.5,0) arc(0:30:1) (0.11,0.14666667);
			 \draw (0.65, 0.2) node {θ};
		\end{tikzpicture}		
´

		\column{0.4\textwidth}
		Γενικότερα θα αναλύσουμε τις γωνίες και τους σχετικούς τριγωνομετρικούς αριθμούς βάσει το πως σχηματίζονται μεταξύ μιας ακτίνας και των οριζόντιων ή κάθετων πλευρών που σχηματίζονται ως \textbf{συνδέσεων των προβολών} με το σημείο στον οριζόντιο άξονα (ή και της υποτείνουσας που προκύπτει από το προϊόν τρίγωνο).

	\end{columns}
\end{frame}

\begin{frame}
	\frametitle{Ασκήσεις Δειγματοληψίας -  Τριγωνομετρία}
	\begin{columns}
		\column{0.4\textwidth}
		
		\begin{tikzpicture}
			\draw (-2,0) -- (2,0);
			\draw (0,2) -- (0,-2);
			\draw (0,0) circle (1);
			\draw (2,0) node[right] {$x$};
			\draw (0.5,0.5) node {\tiny{1o}};
			\draw (-2,0) node[left] {$-x$};
			\draw (-0.5,0.5) node {\tiny{2o}};
			\draw (0,2) node[right] {$y$};
			\draw (-0.5,-0.5) node {\tiny{3o}};
			\draw (0,-2) node[right] {$-y$};
			\draw (0.5,-0.5) node {\tiny{4o}};
		\end{tikzpicture}		
		
		\column{0.6\textwidth}
		\centering
		Χωρίζεται σε τέσσερα τεταρτημόρια ως εξής:
		\begin{itemize}
			\item 1ο Τεταρτημόριο: 0 εως 90 μοίρες ($0$ εως $\frac{\pi}{2}$ ακτίνια)
			\item 2ο Τεταρτημόριο: 90 εως 180 μοίρες ($\frac{\pi}{2}$ εως $\pi$ \textlatin{rad})
			\item 3ο Τεταρτημόριο: 180 εως 270 μοίρες ($\pi$ εως $\frac{3\pi}{2}$ \textlatin{rad})
			\item 4ο Τεταρτημόριο: 270 εως 360 μοίρες ($\frac{3\pi}{2}$ εως $2\pi$ \textlatin{rad})
		 \end{itemize}
	\end{columns}
~\\
Θυμίζουμε πως η μετατροπή μοιρών σε ακτίνια γίνεται με τη σχέση $\alpha = \frac{\pi}{180} * \delta$, όπου $\alpha$ τα ακτίνια και $\delta$ η γωνία σε μοίρες. Η μονάδα μέτρησης είναι $rad$.
\end{frame}

\begin{frame}
	\frametitle{Ασκήσεις Δειγματοληψίας -  Τριγωνομετρία}
	Με τις σχετικές αποδείξεις να παρατίθενται σε αργότερο χρόνο σε χειρόγραφη μορφή για το παρόν \textlatin{thread}, αναφέρουμε εδώ ότι μέσω των σχέσεων που μπορούν να εμφανιστούν εδώ, δύναται να λάβουμε τους τριγωνομετρικούς αριθμούς κάθε γωνίας, αλλά να δούμε και κάποια άλλα πράγματα. Για παράδειγμα:
	\begin{itemize}
		\item Έστω ότι θέλουμε να υπολογίσουμε ένα ημίτονο. Αυτό υπολογίζεται ως το πηλίκο της διαίρεσης της απέναντι της γωνίας πλευράς με την υποτείνουσα του τριγώνου. \pause Αν γνωρίζουμε ότι η υποτείνουσα είναι 1 καθώς είναι η ακτίνα, τότε η απέναντι πλευρά ξέρουμε ότι είναι ίση με το εκάστοτε ημίτονο.
		\item Με παρόμοια λογική λειτουργεί και το συνημίτονο, όπως θα αναδειχθεί και στην επόμενη διαφάνεια.
		\item Διαισθητικά, καταλαβαίνουμε πως και αν προβούμε σε πρόσθεση 2π σε μια γωνία, επιστρέφουμε πάλι σε αυτήν άρα $\cos(x) = \cos(x+2\pi)$ και $\sin(x) = \sin(x+2\pi)$.
	\end{itemize}
\end{frame}

% Prove how sines work
\begin{frame}
	\frametitle{Ασκήσεις Δειγματοληψίας -  Τριγωνομετρία}
	\begin{columns} 
		\column{0.6\textwidth}
		\begin{tikzpicture}
			\draw (-3,0) -- (3,0);
			\draw (0,3) -- (0,-3);
			\draw (0,0) circle (2);
			\draw (3,0) node[right] {$x$};
			\draw (-3,0) node[left] {$-x$};
			\draw (0,3) node[right] {$y$};
			\draw (0,-3) node[right] {$-y$};
			% radius and triangle
			\draw[red] (0,0) -- (0,2);
			%\draw[thick, blue] (1.2,1.6) -- (1.2,0);
			%\draw[thick,orange] (0,0) -- (1.2,0);
			%\draw (1.0,1.5) node[below left] {{\textcolor{red}{$r$}}};
			%\draw (0.6,-0.15) node {{\textcolor{orange}{$x'$}}};
			%\draw (1.20,1.20) node[below right] {{\textcolor{blue}{$y'$}}};
			% \draw (0.5,0) arc(0:90:1) ()
		\end{tikzpicture}		
		
		
		\column{0.4\textwidth}
		\only<1>
		{Υποθέστε πως η γωνία που σχηματίζει μια ευθεία που σχεδιάζουμε από το σημείο Ο(0,0) με τον οριζόντιο άξονα είναι 90 μοίρες (δηλαδή έχουμε προχωρήσει 90 μοίρες από την αρχή του κύκλου).}
		\only<2>
		{Τώρα ας δούμε το εξής:
			\begin{itemize}
				
				\item Αν θέλω να δω που είναι η προβολή του σημείου που καταλήγουμε στην ευθεία $x$ βλέπουμε πως είμαστε ακριβώς επάνω, άρα δεν ορίζεται κάποια ευθεία άρα $|x'|=0$.
			\end{itemize}
		}
		
	\end{columns}
\end{frame}


\begin{frame}
	\frametitle{Ασκήσεις Δειγματοληψίας -  Τριγωνομετρία}
	\begin{columns} 
		\column{0.6\textwidth}
		\begin{tikzpicture}
			\draw (-3,0) -- (3,0);
			\draw (0,3) -- (0,-3);
			\draw (0,0) circle (2);
			\draw (3,0) node[right] {$x$};
			\draw (-3,0) node[left] {$-x$};
			\draw (0,3) node[right] {$y$};
			\draw (0,-3) node[right] {$-y$};
			% radius and triangle
			\draw[red] (0,0) -- (0,2);
			%\draw[thick, blue] (1.2,1.6) -- (1.2,0);
			%\draw[thick,orange] (0,0) -- (1.2,0);
			%\draw (1.0,1.5) node[below left] {{\textcolor{red}{$r$}}};
			%\draw (0.6,-0.15) node {{\textcolor{orange}{$x'$}}};
			%\draw (1.20,1.20) node[below right] {{\textcolor{blue}{$y'$}}};
			% \draw (0.5,0) arc(0:90:1) ()
		\end{tikzpicture}		
		
		
		\column{0.4\textwidth}
		Τώρα ας δούμε το εξής:
			\begin{itemize}
				\item Αν θέλω να δω που είναι η προβολή του σημείου που καταλήγουμε στην ευθεία $y$ βλέπουμε πως θα ταυτίζεται με την ακτίνα μας καθώς θα είναι μια κάθετη ευθεία γραμμή από το σημείο (0,1) στο (0,0) το οποίο κοιτά ακριβώς κάθετα. Άρα $|y'|=1$.
			\end{itemize}
		
		
	\end{columns}
\end{frame}

\begin{frame}
	\frametitle{Ασκήσεις Δειγματοληψίας -  Τριγωνομετρία}
	\begin{columns} 
		\column{0.6\textwidth}
		\begin{tikzpicture}
			\draw (-3,0) -- (3,0);
			\draw (0,3) -- (0,-3);
			\draw (0,0) circle (2);
			\draw (3,0) node[right] {$x$};
			\draw (-3,0) node[left] {$-x$};
			\draw (0,3) node[right] {$y$};
			\draw (0,-3) node[right] {$-y$};
			% radius and triangle
			\draw[red] (0,0) -- (0,2);
			%\draw[thick, blue] (1.2,1.6) -- (1.2,0);
			%\draw[thick,orange] (0,0) -- (1.2,0);
			%\draw (1.0,1.5) node[below left] {{\textcolor{red}{$r$}}};
			%\draw (0.6,-0.15) node {{\textcolor{orange}{$x'$}}};
			%\draw (1.20,1.20) node[below right] {{\textcolor{blue}{$y'$}}};
			% \draw (0.5,0) arc(0:90:1) ()
		\end{tikzpicture}		
		
		
		\column{0.4\textwidth}
		Όταν αναφερόμαστε σε σχήματα αναπαριστώμενα επί του Καρτεσιανού συστήματος συντεταγμένων, οι τύποι υπολογισμού των συνημίτονων και ημίτονων είναι:
		\begin{itemize}
			\item συνθ = $\cos(\theta) = \frac{x'}{r}$
			\item ημθ = $\sin(\theta) = \frac{y'}{r}$
		\end{itemize}
	\end{columns}
, όπου $r$ η ακτίνα και $x'$ και $y'$ τα ευθύγραμμα τμήματα που αναφέραμε πριν (αυτά που ενώνουν τις αντίστοιχες προβολές με τους άξονες).
\end{frame}


\begin{frame}
	\frametitle{Ασκήσεις Δειγματοληψίας -  Τριγωνομετρία}
	\begin{columns} 
		\column{0.6\textwidth}
		\begin{tikzpicture}
			\draw (-3,0) -- (3,0);
			\draw (0,3) -- (0,-3);
			\draw (0,0) circle (2);
			\draw (3,0) node[right] {$x$};
			\draw (-3,0) node[left] {$-x$};
			\draw (0,3) node[right] {$y$};
			\draw (0,-3) node[right] {$-y$};
			% radius and triangle
			\draw[red] (0,0) -- (0,2);
			%\draw[thick, blue] (1.2,1.6) -- (1.2,0);
			%\draw[thick,orange] (0,0) -- (1.2,0);
			%\draw (1.0,1.5) node[below left] {{\textcolor{red}{$r$}}};
			%\draw (0.6,-0.15) node {{\textcolor{orange}{$x'$}}};
			%\draw (1.20,1.20) node[below right] {{\textcolor{blue}{$y'$}}};
			% \draw (0.5,0) arc(0:90:1) ()
		\end{tikzpicture}		
		
		
		\column{0.4\textwidth}
		Έχουμε λοιπόν:
		\begin{itemize}
			\item $\cos(\frac{\pi}{2}) = \frac{x'}{r} = \frac{0}{1} = 0$
			\item $\sin(\frac{\pi}{2}) = \frac{y'}{r} = \frac{1}{1} = 1$
		\end{itemize}
	Άρα υπολογίσαμε πως το ημίτονο των 90 μοιρών είναι 1, ενώ το συνημίτονο 0. Με παρόμοιο τρόπο μπορούμε να υπολογίσουμε και άλλους τριγωνομετρικούς αριθμούς (βρίσκετε ονλάιν κάποιο πίνακα σχετικό).
	\end{columns}
\end{frame}

\begin{frame}
	% Please add the following required packages to your document preamble:
	% \usepackage[table,xcdraw]{xcolor}
	% Beamer presentation requires \usepackage{colortbl} instead of \usepackage[table,xcdraw]{xcolor}
	\begin{table}[]
		\centering
		\begin{tabular}{|l|llllllll|}
			\hline
			& \multicolumn{8}{c|}{\cellcolor[HTML]{AB5E5E}Γωνία (σε ακτίνια)}                                                                                                                                                                                                \\ \hline
			\multicolumn{1}{|c|}{\cellcolor[HTML]{AB5E5E}Τριγωνομετρικοί Αριθμοί} & \multicolumn{1}{l|}{0} & \multicolumn{1}{l|}{$\frac{\pi}{6}$}      & \multicolumn{1}{l|}{$\frac{\pi}{4}$}      & \multicolumn{1}{l|}{$\frac{\pi}{3}$}      & \multicolumn{1}{l|}{$\frac{\pi}{2}$} & \multicolumn{1}{l|}{$\pi$} & \multicolumn{1}{l|}{$\frac{3\pi}{2}$} & 2$\pi$ \\ \hline
			$\sin(\theta)$                                                        & \multicolumn{1}{l|}{0} & \multicolumn{1}{l|}{$\frac{1}{2}$}        & \multicolumn{1}{l|}{$\frac{\sqrt{2}}{2}$} & \multicolumn{1}{l|}{$\frac{\sqrt{3}}{2}$} & \multicolumn{1}{l|}{1}               & \multicolumn{1}{l|}{0}     & \multicolumn{1}{l|}{-1}               & 0      \\ \hline
			$\cos(\theta)$                                                        & \multicolumn{1}{l|}{1} & \multicolumn{1}{l|}{$\frac{\sqrt{3}}{2}$} & \multicolumn{1}{l|}{$\frac{\sqrt{2}}{2}$} & \multicolumn{1}{l|}{$\frac{1}{2}$}        & \multicolumn{1}{l|}{0}               & \multicolumn{1}{l|}{-1}    & \multicolumn{1}{l|}{0}                & 1      \\ \hline
		\end{tabular}
	\end{table}
\end{frame}

\begin{frame}
	\frametitle{Ασκήσεις Δειγματοληψίας - Τριγωνομετρία (συν.)}
		\begin{columns} 
		\column{0.6\textwidth}
		\begin{tikzpicture}
			\draw (-3,0) -- (3,0);
			\draw (0,3) -- (0,-3);
			\draw (0,0) circle (2);
			\draw (3,0) node[right] {$x$};
			\draw (-3,0) node[left] {$-x$};
			\draw (0,3) node[right] {$y$};
			\draw (0,-3) node[right] {$-y$};
			% radius and triangle
			\draw[red] (0,0) -- (1.2,-1.6);
			\draw[blue] (0,0) -- (1.2,1.6);
			%\draw[thick, blue] (1.2,1.6) -- (1.2,0);
			%\draw[thick,orange] (0,0) -- (1.2,0);
			%\draw (1.0,1.5) node[below left] {{\textcolor{red}{$r$}}};
			%\draw (0.6,-0.15) node {{\textcolor{orange}{$x'$}}};
			%\draw (1.20,1.20) node[below right] {{\textcolor{blue}{$y'$}}};
			% \draw (0.5,0) arc(0:90:1) ()
		\end{tikzpicture}		
		
		
		\column{0.4\textwidth}
		\begin{itemize}
		\item	Έστω ότι η μπλε ευθεία δείχνει στις 45 μοίρες ($\frac{\pi}{4}$ ακτίνια) και η κόκκινη στα $-\frac{\pi}{4}$.
		\item Αν προσέξουμε, θα δούμε ότι στον οριζόντιο άξονα έχουν την ίδια ακριβώς προβολή! Επομένως, $|x'| = \frac{\sqrt{2}}{2}$
	\end{itemize}
	\end{columns}
\end{frame}

\begin{frame}
	\frametitle{Ασκήσεις Δειγματοληψίας - Τριγωνομετρία (συν.)}
	\begin{columns} 
		\column{0.6\textwidth}
		\begin{tikzpicture}
			\draw (-3,0) -- (3,0);
			\draw (0,3) -- (0,-3);
			\draw (0,0) circle (2);
			\draw (3,0) node[right] {$x$};
			\draw (-3,0) node[left] {$-x$};
			\draw (0,3) node[right] {$y$};
			\draw (0,-3) node[right] {$-y$};
			% radius and triangle
			\draw[red] (0,0) -- (1.2,-1.6);
			\draw[blue] (0,0) -- (1.2,1.6);
			%\draw[thick, blue] (1.2,1.6) -- (1.2,0);
			%\draw[thick,orange] (0,0) -- (1.2,0);
			%\draw (1.0,1.5) node[below left] {{\textcolor{red}{$r$}}};
			%\draw (0.6,-0.15) node {{\textcolor{orange}{$x'$}}};
			%\draw (1.20,1.20) node[below right] {{\textcolor{blue}{$y'$}}};
			% \draw (0.5,0) arc(0:90:1) ()
		\end{tikzpicture}		
		
		
		\column{0.4\textwidth}
		\begin{itemize}
			\item Εφαρμόζοντας τον τύπο του συνημιτόνου έχουμε $\cos(-\frac{\pi}{4}) = \frac{x'}{r} = \cos(\frac{\pi}{4})$. Γενικά λοιπόν, έχουμε πως \alert{$\cos(-x) = \cos(x)$}.
			\item Επίσης, σημειώνουμε πως όπως είχαμε $\sin(x+2\pi) = \sin(x)$ και $\cos(x+2\pi) = \cos(x)$, μπορεί να έχουμε αμέτρητες περιστροφές, άρα ισχύει πως \alert{$\sin(x+k*2\pi) = \sin(x)$} και \alert{$\cos(x+k*2\pi) = \cos(x)$}, με $k \in \mathbb{Z}$.
		\end{itemize}
	\end{columns}
\end{frame}

\begin{frame}
	\frametitle{Ασκήσεις Δειγματοληψίας - 1η (Λύση-Εξήγηση)}
		Έχουμε λοιπόν $x[n] = 2\cos(\frac{4\pi n}{3})$. Για να μετασχηματιστεί σε μια απλούστερη μορφή, προχωρούμε ως εξής:
		\begin{itemize}
			\item Θέλουμε να εμφανίσουμε πρώτα ένα πολλαπλάσιο ακέραιο του 2π, ώστε να μπορούμε να εφαρμόσουμε την ιδιότητα $\cos(x+2\pi) = \cos(x)$. Εδώ, μπορούμε να πούμε πως $x[n] = 2\cos(\frac{4\pi n}{3}) \rightarrow x[n] = 2\cos(\frac{6\pi n}{3} - \frac{2\pi n}{3}) \rightarrow x[n] = 2\cos(2\pi n - \frac{2\pi n}{3}) $. Έχουμε $n$ περιστροφές του κύκλου άρα καταλήγουμε σε $x[n] = 2\cos(\cancel{2\pi n} - \frac{2\pi n}{3}) \rightarrow 2\cos(- \frac{2\pi n}{3})$
			\pause
			\item Καθώς δείξαμε ότι $\cos(-x) = \cos(x)$, έχουμε μετ''ιδίας λογικής $x[n] = 2\cos(- \frac{2\pi n}{3}) \rightarrow \text{\textcolor{red}{$x[n] = 2\cos(\frac{2\pi n}{3})$}} $
		\end{itemize}
\end{frame}

\begin{frame}
	\frametitle{Ασκήσεις Δειγματοληψίας - 1η (Λύση-Εξήγηση)}
	Για να χτίσουμε τώρα το σήμα συνεχούς χρόνου που προκύπτει από την ανακατασκευή, θα λάβουμε τον τύπο
	\begin{equation*}
		f = F * T_s
	\end{equation*}
	και θα τον μετασχηματίσουμε σε
	\begin{equation*}
		f = F * T_s \rightarrow F = f * \frac{1}{T_s} \rightarrow F = f*F_s
	\end{equation*}

\end{frame}

\begin{frame}
	\frametitle{Ασκήσεις Δειγματοληψίας - 1η (Λύση-Εξήγηση)}
	Το $x[n]$ μας παρατηρούμε πως είναι της μορφής $A\cos[2\pi f n]$ , συγκεκριμένα δε
	\begin{align*}
		\omega = 2\pi f \rightarrow \frac{2\pi}{3} = 2\pi f \\ 
		\rightarrow  \frac{2\pi}{3} * \frac{1}{2\pi} = 2\pi f * \frac{1}{2\pi} \rightarrow f = \frac{1}{3}
	\end{align*} 
	Συνεχίζοντας, έχουμε τον τύπο υπολογισμού της διακριτής συχνότητας με τη βοήθεια της αναλογικής και της περιόδου δειγματοληψίας. Αντικαθιστώντας, έχουμε:
	\begin{equation*}
		f = F' * T_s \rightarrow F' = f * \frac{1}{T_s} \rightarrow F' = f*F_s \rightarrow F' = \frac{1}{3} * 75
	\end{equation*}
	Άρα $F' = 25 Hz$.
	
\end{frame}

\begin{frame}
	\frametitle{Ασκήσεις Δειγματοληψίας - 1η (Λύση-Εξήγηση)}
	Άρα $F' = 25 Hz$ και αν τοποθετηθεί στο ως συχνότητα αναλογικού ημιτονοειδούς σήματος ίδιου πλάτους, έχουμε $x'(t) = 2\cos(2\pi25t)$ ή $x'(t) = 2\cos(50\pi t)$.\\~\\
	
	Υπάρχουν ψευδείς συχνότητες? \textbf{Ναι!} Μια εξ αυτών είναι η $F' = 25Hz$. \\~\\
	
	Αυτή η ψευδής συχνότητα, όπως αναφέραμε, είναι προϊόν της ασυμφωνίας με το Θεώρημα Δειγματοληψίας (ή Θεώρημα \textlatin{Nyquist-Shannon}).
	
\end{frame}

\begin{frame}
	\frametitle{Ασκήσεις Δειγματοληψίας - Οι Υπόλοιπες?}
	Για υπόλοιπες ασκήσεις δειγματοληψίας (και αν προλάβουμε κάποια πραγματάκια σε \textlatin{Matlab}), θα αναφερθείτε στα υπόλοιπα αρχεία που θα συντελούν το ενισχυτικό υλικό του \textlatin{thread}. Αν υπάρχουν απορίες ή προτάσεις για διορθώσεις, αναφέρετε τις στο κανάλι του μαθήματος. \\~\\
	\begin{center}
		Καλή συνέχεια, και καλά διαβάσματα!
	\end{center}
	
\end{frame}
\section{Βιβλιογραφία}

		\begin{frame}
		\frametitle{Πηγές}
\bibliographystyle{custard}

\bibliography{fs}   % references.bib, without the extension
\nocite{*}
		\end{frame}
\end{document}